\section*{Chapitre \ref{ch:rs:vectorised-backtests} --- Backtests vectorisés}

\begin{solution}{exo:rs:vectorised-backtests:1}
Le rendement brut vaut $12\%/252 = 0.0476\%$ par jour~; d'après la \cref{prop:rs:vectorised-backtests:breakeven}, le coût de point mort est
$0.0476~\%/(2 \times 0.2) = 0.119~\%$, soit 11.9 points de base par unité échangée.
\end{solution}

\begin{solution}{exo:rs:vectorised-backtests:2}
Au terme de la période, A vaut $0.66$ et B $0.40$ d'un capital de 1.06~: poids de 62.3~\% et 37.7~\%. Pour revenir à 60/40, vendre 2.3~\% du capital en A et acheter
2.3~\% en B.
\end{solution}

\begin{solution}{exo:rs:vectorised-backtests:3}
Par composition~: $(1.5 \times 0.6)^5 = 0.9^5 = 0.59$, soit une perte de 41~\%. Par addition~: $5 \times 50~\% - 5 \times 40~\% = +50~\%$, un gain qui ne s'est jamais produit.
\end{solution}

\begin{solution}{exo:rs:vectorised-backtests:4}
Le signal du retournement est l'opposé du rendement du jour, donc le multiplier par ce rendement donne l'opposé du carré de celui-ci~: une perte certaine. Un signal de momentum qui
intègre le rendement du jour donne un multiple positif de ce carré~: un gain certain. Le look-ahead ajoute la corrélation du signal avec le rendement à partir duquel il a été
construit, quel que soit son signe.
\end{solution}

\begin{solution}{exo:rs:vectorised-backtests:5}
Un univers de survivants écarte les valeurs radiées par la suite, et avec elles leurs rendements avant la radiation et lors de celle-ci. Une stratégie qui aurait détenu ces
valeurs en position longue voit disparaître leurs pertes (son \omterm{def:rs:vectorised-backtests:backtest}{backtest} monte)~; une stratégie qui les aurait détenues à découvert voit disparaître leurs gains (son \omterm{def:rs:vectorised-backtests:backtest}{backtest} baisse). Dans le marché synthétique, les
valeurs radiées étaient les perdants passés du momentum, détenues à découvert, et leurs baisses jusqu'à la radiation rapportaient des profits à la jambe vendeuse.
\end{solution}

\begin{solution}{exo:rs:vectorised-backtests:6}
Un portefeuille neutre en dollars a une exposition nette nulle, donc son cash est son capital~: $+3~\%$ par an d'intérêts. Ses positions vendeuses représentent 1 (la moitié de l'exposition brute de 2) et coûtent $0.5~\% \times 1
= 0.5~\%$. Au total~: $+2.5~\%$ par an. (Le ratio de Sharpe se calcule sur le rendement excédentaire par rapport au taux monétaire.)
\end{solution}

\begin{solution}{exo:rs:vectorised-backtests:7}
\texttt{rs\_vecbt.smoothed(3)}~: la \omterm{def:rs:vectorised-backtests:turnover}{rotation du portefeuille} (turnover) tombe de 75~\% à 42~\% du portefeuille par jour et le ratio de Sharpe brut de 6.95 à 4.13, tandis que le coût de point mort
varie à peine (7.4 points de base contre 7.3)~: le lissage retire de l'alpha en même temps que les transactions, car l'information du retournement est vieille d'un jour. Il n'aide
que les signaux dont l'information dure plus longtemps que leur bruit.
\end{solution}

\begin{solution}{exo:rs:vectorised-backtests:8}
Les ordres au marché à la clôture (MOC) sont transmis avant la clôture, généralement quelques minutes avant, alors que le cours de clôture n'est pas connu~; un signal calculé à partir de ce cours n'est
pas disponible au moment où l'ordre doit être transmis. Soit calculer le signal à partir de cours connus avant l'heure limite de transmission, soit exécuter à la clôture suivante~; et modéliser
les coûts propres à l'enchère de clôture.
\end{solution}

\begin{solution}{pb:rs:vectorised-backtests:1}
\begin{enumerate}
\item Ratio de Sharpe 7.17, 27.0~\% par an, en supposant l'exécution à la clôture même où le signal a été calculé, sur les survivants d'aujourd'hui, sans coûts.
\item $-65$, soit $-694~\%$ par an~: le signal est l'opposé du rendement du jour, multiplié par ce rendement.
\item 74~\% du portefeuille par jour.
\item Les rendements jusqu'à la clôture, mais pas la possibilité d'échanger à cette clôture, ni la liste des survivants.
\item Retournement~: 7.22 sur l'univers connu à chaque date, 6.95 sur les 500 titres les plus liquides. Momentum~: 0.44, puis 0.32.
\item Les titres radiés étaient les perdants du momentum, détenus à découvert~; les écarter supprime les profits de la jambe vendeuse (0.21 contre 0.44).
\item Il faut le relever~: les titres en difficulté ayant fait faillite feraient défaut aux positions longues du \omterm{def:rs:vectorised-backtests:backtest}{backtest}.
\item Retournement~: de 6.95 à $-2.58$. Momentum~: de 0.32 à 0.21.
\item 7.3 points de base par unité échangée.
\item Environ 0.25~\% par an pour le retournement~: de $-2.58$ à $-2.64$.
\item $-9.45$~: l'information du retournement ne dure qu'un jour, si bien qu'un portefeuille échangé un jour en retard ne rapporte aucun alpha et paie tout de même ses coûts.
\item Par composition, le capital tombe à 0.35 de son niveau initial~; par addition, le cumul atteint $-0.04$.
\item \textbf{Résultat nommé.} Retournement~: 7.17 (naïf), 7.22 (univers connu à chaque date), 6.95 (titres liquides), $-2.58$ (coûts), $-2.64$ (frais d'emprunt), $-9.45$ (clôture suivante).
Momentum~: 0.21, 0.44, 0.32, 0.21, 0.19, 0.17. Le retournement s'effondre sous les coûts, puis de nouveau sous le \omterm{def:rs:vectorised-backtests:timing}{délai d'exécution}~; le momentum perd 60~\% de son meilleur chiffre mais conserve un
ratio de Sharpe positif et faible à travers toutes les corrections.
\item Le momentum est un résultat de niveau 1~: lent, liquide, insensible à l'exécution. Le retournement exige que l'enchère de clôture soit modélisée (niveau~2 avec un modèle
d'exécution à l'enchère, ou niveau~3), et sa réponse de niveau 1 est déjà négative.
\item Cela donnerait au retournement un moyen légitime d'échanger à la clôture~; l'alpha qui subsiste dépend de la part du mouvement de la journée qui intervient après
15:50, et le coût de l'enchère remplace les 10 points de base génériques.
\item N'échanger que les titres dont le signal franchit un seuil, conserver les positions jusqu'à ce que le signal s'inverse (un spread achat-conservation, comme le recommandent Novy-Marx et Velikov),
ou combiner le signal avec des signaux plus lents~; chacune de ces règles réduit davantage la rotation que l'alpha, seulement si l'alpha n'est pas tout entier le premier jour.
\item Les deux premières lignes~: le \omterm{def:rs:vectorised-backtests:timing}{délai d'exécution} et l'univers connu à chaque date.
\item Qu'un retournement quotidien, dont la rotation dépasse largement 50~\% par mois, ne résiste pas aux coûts, sauf si ceux-ci sont exceptionnellement faibles.
\item Chaque variante testée, avec ses réglages (univers, délai, coûts), la cascade de niveau 1 et la décision prise (chapitre~1).
\item Trier les idées rapidement et honnêtement lorsque l'exécution ne change pas la réponse.
\end{enumerate}
\end{solution}

\begin{solution}{iq:rs:vectorised-backtests:1}
Look-ahead (entrées non connues à l'\omterm{def:rs:vectorised-backtests:timing}{instant de décision}, ou exécution au cours même utilisé par le signal), \omterm{def:rs:point-in-time-data-and-the-biases:survivorship}{biais de survie} (univers d'aujourd'hui), coûts nuls ou faibles, absence de frais d'emprunt
ou de financement, et data snooping (nombreuses variantes, une seule rapportée)~; aussi le cumul arithmétique de grands rendements, et l'oubli de la capacité.
\end{solution}

\begin{solution}{iq:rs:vectorised-backtests:2}
Construire le panel de signaux avec chaque donnée à son \omterm{def:rs:point-in-time-data-and-the-biases:bitemporal}{temps de connaissance}~; le filtrer sur l'\omterm{def:rs:universe-symbology-and-corporate-actions:universe}{univers investissable} à chaque date~; le convertir en poids cibles~; les décaler à la
première période où ils peuvent être détenus~; calculer le portefeuille après dérive des cours et les transactions vers la nouvelle cible~; imputer les coûts aux transactions, les frais d'emprunt aux positions vendeuses, le financement du cash et
de l'effet de levier~; composer les rendements. Les biais interviennent à chaque étape~: délai (décalage), univers (survivants), coûts, financement, composition.
\end{solution}

\begin{solution}{iq:rs:vectorised-backtests:3}
Sa rotation peut faire que les coûts dépassent son rendement brut (le coût de point mort peut être inférieur aux coûts réalistes)~; elle peut dépendre d'une exécution à un prix
qu'elle ne pouvait pas obtenir~; elle peut ne pas passer à l'échelle (capacité) ou être l'une des nombreuses variantes testées. Sur le retournement synthétique, le ratio est passé de 7 à $-2.6$ avec 10 points de base de coût.
\end{solution}

\begin{solution}{iq:rs:vectorised-backtests:4}
Suffisant pour des stratégies lentes et liquides dont les transactions restent petites par rapport aux volumes et dont les coûts se résument à une fourchette par unité échangée. Un \omterm{def:rs:vectorised-backtests:backtest}{backtest} événementiel (ou par rejeu) est
nécessaire lorsque les exécutions dépendent du carnet d'ordres (ordres passifs ou ordres à cours limité), lorsque le calage intrajournalier ou les enchères comptent, lorsque l'impact propre de la stratégie compte,
ou lorsque des événements interviennent entre la décision et la transaction.
\end{solution}

\begin{solution}{iq:rs:vectorised-backtests:5}
Imputer les frais d'emprunt titre par titre aux positions vendeuses détenues à chaque période (avec des limites de disponibilité), créditer les intérêts sur le cash et sur le produit des ventes à découvert, comme le prime broker
le verse, prélever le spread de financement sur les liquidités empruntées pour l'effet de levier, et calculer le ratio de Sharpe sur le rendement excédentaire par rapport au taux monétaire.
\end{solution}

\begin{solution}{iq:rs:vectorised-backtests:6}
Par période, le portefeuille échange $2\tau$ de poids et paie $2\tau c$~; le rendement net vaut $g - 2\tau c$, nul pour $c = g/(2\tau)$. Pour le retournement synthétique, $g$ vaut environ
0.109~\% par jour et $\tau = 0.75$, soit 7.3 points de base~; pour le momentum mensuel, $g$ vaut environ 0.011~\% par jour et $\tau$ vaut environ 0.019, soit environ 30 points de base~: quatre fois la
tolérance aux coûts du retournement, avec un dixième de son rendement brut.
\end{solution}
